% Authored for a11y-stem's development set (dev-v2). Original text, no third-party material.
% Licence: MIT. Author: a11y-stem contributors.
\documentclass[11pt]{article}
% Computer Modern text fonts, as pdfLaTeX writes them; XeTeX defaults to Latin Modern.
\usepackage[OT1]{fontenc}
\renewcommand{\rmdefault}{cmr}
\renewcommand{\ttdefault}{cmtt}
\usepackage{amsmath,amssymb}
\usepackage[margin=1in]{geometry}
\setlength{\parindent}{0pt}\setlength{\parskip}{8pt}
\begin{document}
\begin{center}{\Large\bfseries Atoms, Ions and Reactions: Notation for Chemistry and Nuclear Science}\end{center}
\section*{Formulas and Ions}
Chemical formulas are set upright, with subscripts for the number of atoms and superscripts for charge.
The formula of water is \label{eq000}$\mathrm{H_2O}$. The formula of calcium hydroxide is \label{eq001}$\mathrm{Ca(OH)_2}$. The formula of glucose is \label{eq002}$\mathrm{C_6H_{12}O_6}$. The formula of ammonia is \label{eq003}$\mathrm{NH_3}$. The formula of sulfuric acid is \label{eq004}$\mathrm{H_2SO_4}$.
The formula of ethanol is \label{eq005}$\mathrm{C_2H_6O}$. The formula of aluminium sulfate is \label{eq006}$\mathrm{Al_2(SO_4)_3}$. The formula of sodium bicarbonate is \label{eq007}$\mathrm{NaHCO_3}$. The formula of methane is \label{eq008}$\mathrm{CH_4}$. The formula of iron(III) oxide is \label{eq009}$\mathrm{Fe_2O_3}$.
The sodium cation is written \label{eq010}$\mathrm{Na^+}$ and carries a charge of 1. The chloride anion is written \label{eq011}$\mathrm{Cl^-}$ and carries a charge of 1. The calcium cation is written \label{eq012}$\mathrm{Ca^{2+}}$ and carries a charge of 2. The sulfate anion is written \label{eq013}$\mathrm{SO_4^{2-}}$ and carries a charge of 2. The iron(III) cation is written \label{eq014}$\mathrm{Fe^{3+}}$ and carries a charge of 3.
The ammonium cation is written \label{eq015}$\mathrm{NH_4^+}$ and carries a charge of 1. The hydroxide anion is written \label{eq016}$\mathrm{OH^-}$ and carries a charge of 1. The phosphate anion is written \label{eq017}$\mathrm{PO_4^{3-}}$ and carries a charge of 3. The copper(II) cation is written \label{eq018}$\mathrm{Cu^{2+}}$ and carries a charge of 2. The nitrate anion is written \label{eq019}$\mathrm{NO_3^-}$ and carries a charge of 1.
Aluminium sulfate, Al\textsubscript{2}(SO\textsubscript{4})\textsubscript{3}, is made of molecules or ions that differ from those of ammonia, NH\textsubscript{3}. The copper(II) ion, Cu$^{2+}$, is formed when the neutral species loses 2 electrons. When the salt of calcium and phosphate dissolves in water it releases Ca$^{2+}$ and PO\textsubscript{4}$^{3-}$ ions. A solution of methane (CH\textsubscript{4}) was mixed with a solution containing NH\textsubscript{4}$^{+}$ ions.

In the reaction NH\textsubscript{3} reacts with H\textsubscript{2}SO\textsubscript{4} to give H\textsubscript{2}O, the mass of the products equals the mass of the reactants. Water has the formula H\textsubscript{2}O and is stored in a sealed bottle. A bottle labelled chloride holds the ion Cl$^{-}$, which is surrounded by water molecules. A bottle labelled nitrate holds the ion NO\textsubscript{3}$^{-}$, which is surrounded by water molecules.

In the reaction C\textsubscript{2}H\textsubscript{6}O reacts with Al\textsubscript{2}(SO\textsubscript{4})\textsubscript{3} to give CH\textsubscript{4}, the mass of the products equals the mass of the reactants. Sodium bicarbonate, NaHCO\textsubscript{3}, is made of molecules or ions that differ from those of sulfuric acid, H\textsubscript{2}SO\textsubscript{4}. A solution of aluminium sulfate (Al\textsubscript{2}(SO\textsubscript{4})\textsubscript{3}) was mixed with a solution containing Fe$^{3+}$ ions. When the salt of sodium and phosphate dissolves in water it releases Na$^{+}$ and PO\textsubscript{4}$^{3-}$ ions.

Glucose has the formula C\textsubscript{6}H\textsubscript{12}O\textsubscript{6} and is stored in a sealed bottle. The calcium ion, Ca$^{2+}$, is formed when the neutral species loses 2 electrons. A solution of aluminium sulfate (Al\textsubscript{2}(SO\textsubscript{4})\textsubscript{3}) was mixed with a solution containing Fe$^{3+}$ ions. In the reaction Fe\textsubscript{2}O\textsubscript{3} reacts with C\textsubscript{6}H\textsubscript{12}O\textsubscript{6} to give Ca(OH)\textsubscript{2}, the mass of the products equals the mass of the reactants.

When the salt of copper(II) and sulfate dissolves in water it releases Cu$^{2+}$ and SO\textsubscript{4}$^{2-}$ ions. The ammonium ion, NH\textsubscript{4}$^{+}$, is formed when the neutral species loses one electron. Ammonia has the formula NH\textsubscript{3} and is stored in a sealed bottle. A bottle labelled phosphate holds the ion PO\textsubscript{4}$^{3-}$, which is surrounded by water molecules.

Aluminium sulfate, Al\textsubscript{2}(SO\textsubscript{4})\textsubscript{3}, is made of molecules or ions that differ from those of ethanol, C\textsubscript{2}H\textsubscript{6}O. A solution of sodium bicarbonate (NaHCO\textsubscript{3}) was mixed with a solution containing Cu$^{2+}$ ions. The calcium ion, Ca$^{2+}$, is formed when the neutral species loses 2 electrons. In the reaction Fe\textsubscript{2}O\textsubscript{3} reacts with H\textsubscript{2}SO\textsubscript{4} to give NaHCO\textsubscript{3}, the mass of the products equals the mass of the reactants.

A bottle labelled nitrate holds the ion NO\textsubscript{3}$^{-}$, which is surrounded by water molecules. Iron(III) oxide, Fe\textsubscript{2}O\textsubscript{3}, is made of molecules or ions that differ from those of ammonia, NH\textsubscript{3}. When the salt of copper(II) and nitrate dissolves in water it releases Cu$^{2+}$ and NO\textsubscript{3}$^{-}$ ions. Aluminium sulfate has the formula Al\textsubscript{2}(SO\textsubscript{4})\textsubscript{3} and is stored in a sealed bottle.

A bottle labelled nitrate holds the ion NO\textsubscript{3}$^{-}$, which is surrounded by water molecules. A solution of sodium bicarbonate (NaHCO\textsubscript{3}) was mixed with a solution containing Fe$^{3+}$ ions. The ammonium ion, NH\textsubscript{4}$^{+}$, is formed when the neutral species loses one electron. In the reaction CH\textsubscript{4} reacts with C\textsubscript{2}H\textsubscript{6}O to give Fe\textsubscript{2}O\textsubscript{3}, the mass of the products equals the mass of the reactants.
\section*{Chemical Equations}
Each equation is balanced for atoms and, where ions appear, for charge.
Mixing two clear solutions can form a solid. The net ionic equations below leave out the spectator ions:\nopagebreak
\label{eq020}\[
  \mathrm{Ag^+(aq)} + \mathrm{Cl^-(aq)} \longrightarrow \mathrm{AgCl(s)}
\]
\label{eq021}\[
  \mathrm{Pb^{2+}(aq)} + 2\,\mathrm{I^-(aq)} \longrightarrow \mathrm{PbI_2(s)}
\]
\label{eq022}\[
  \mathrm{Ba^{2+}(aq)} + \mathrm{SO_4^{2-}(aq)} \longrightarrow \mathrm{BaSO_4(s)}
\]
\label{eq023}\[
  \mathrm{Ca^{2+}(aq)} + \mathrm{CO_3^{2-}(aq)} \longrightarrow \mathrm{CaCO_3(s)}
\]
\label{eq024}\[
  \mathrm{Fe^{3+}(aq)} + 3\,\mathrm{OH^-(aq)} \longrightarrow \mathrm{Fe(OH)_3(s)}
\]
\label{eq025}\[
  \mathrm{Cu^{2+}(aq)} + \mathrm{S^{2-}(aq)} \longrightarrow \mathrm{CuS(s)}
\]
\label{eq026}\[
  2\,\mathrm{Ag^+(aq)} + \mathrm{CrO_4^{2-}(aq)} \longrightarrow \mathrm{Ag_2CrO_4(s)}
\]
\label{eq027}\[
  \mathrm{Mg^{2+}(aq)} + 2\,\mathrm{OH^-(aq)} \longrightarrow \mathrm{Mg(OH)_2(s)}
\]
\label{eq028}\[
  3\,\mathrm{Ca^{2+}(aq)} + 2\,\mathrm{PO_4^{3-}(aq)} \longrightarrow \mathrm{Ca_3(PO_4)_2(s)}
\]
An acid and a base give a salt and water. The coefficients are chosen so that every atom is conserved:\nopagebreak
\label{eq029}\[
  \mathrm{HCl(aq)} + \mathrm{NaOH(aq)} \longrightarrow \mathrm{NaCl(aq)} + \mathrm{H_2O(l)}
\]
\label{eq030}\[
  \mathrm{H_2SO_4(aq)} + 2\,\mathrm{KOH(aq)} \longrightarrow \mathrm{K_2SO_4(aq)} + 2\,\mathrm{H_2O(l)}
\]
\label{eq031}\[
  2\,\mathrm{HNO_3(aq)} + \mathrm{Ca(OH)_2(aq)} \longrightarrow \mathrm{Ca(NO_3)_2(aq)} + 2\,\mathrm{H_2O(l)}
\]
\label{eq032}\[
  2\,\mathrm{H_3PO_4(aq)} + 3\,\mathrm{Ba(OH)_2(aq)} \longrightarrow \mathrm{Ba_3(PO_4)_2(s)} + 6\,\mathrm{H_2O(l)}
\]
\label{eq033}\[
  3\,\mathrm{HCl(aq)} + \mathrm{Al(OH)_3(aq)} \longrightarrow \mathrm{AlCl_3(aq)} + 3\,\mathrm{H_2O(l)}
\]
Complete combustion of a fuel in oxygen produces carbon dioxide and water:\nopagebreak
\label{eq034}\[
  \mathrm{CH_4(g)} + 2\,\mathrm{O_2(g)} \longrightarrow \mathrm{CO_2(g)} + 2\,\mathrm{H_2O(l)}
\]
\label{eq035}\[
  \mathrm{C_3H_8(g)} + 5\,\mathrm{O_2(g)} \longrightarrow 3\,\mathrm{CO_2(g)} + 4\,\mathrm{H_2O(l)}
\]
\label{eq036}\[
  \mathrm{C_5H_{12}(l)} + 8\,\mathrm{O_2(g)} \longrightarrow 5\,\mathrm{CO_2(g)} + 6\,\mathrm{H_2O(l)}
\]
\label{eq037}\[
  \mathrm{C_2H_4(g)} + 3\,\mathrm{O_2(g)} \longrightarrow 2\,\mathrm{CO_2(g)} + 2\,\mathrm{H_2O(l)}
\]
\label{eq038}\[
  \mathrm{C_2H_6O(l)} + 3\,\mathrm{O_2(g)} \longrightarrow 2\,\mathrm{CO_2(g)} + 3\,\mathrm{H_2O(l)}
\]
\label{eq039}\[
  \mathrm{C_6H_{12}O_6(s)} + 6\,\mathrm{O_2(g)} \longrightarrow 6\,\mathrm{CO_2(g)} + 6\,\mathrm{H_2O(l)}
\]
Active metals dissolve in acids and release hydrogen gas:\nopagebreak
\label{eq040}\[
  \mathrm{Mg(s)} + 2\,\mathrm{HCl(aq)} \longrightarrow \mathrm{MgCl_2(aq)} + \mathrm{H_2(g)}
\]
\label{eq041}\[
  \mathrm{Zn(s)} + 2\,\mathrm{HCl(aq)} \longrightarrow \mathrm{ZnCl_2(aq)} + \mathrm{H_2(g)}
\]
\label{eq042}\[
  2\,\mathrm{Al(s)} + 3\,\mathrm{H_2SO_4(aq)} \longrightarrow \mathrm{Al_2(SO_4)_3(aq)} + 3\,\mathrm{H_2(g)}
\]
\label{eq043}\[
  \mathrm{Fe(s)} + 2\,\mathrm{HCl(aq)} \longrightarrow \mathrm{FeCl_2(aq)} + \mathrm{H_2(g)}
\]
A more reactive metal displaces a less reactive one from solution:\nopagebreak
\label{eq044}\[
  \mathrm{Zn(s)} + \mathrm{Cu^{2+}(aq)} \longrightarrow \mathrm{Zn^{2+}(aq)} + \mathrm{Cu(s)}
\]
\label{eq045}\[
  \mathrm{Mg(s)} + \mathrm{Fe^{2+}(aq)} \longrightarrow \mathrm{Mg^{2+}(aq)} + \mathrm{Fe(s)}
\]
\label{eq046}\[
  2\,\mathrm{Al(s)} + 3\,\mathrm{Cu^{2+}(aq)} \longrightarrow 2\,\mathrm{Al^{3+}(aq)} + 3\,\mathrm{Cu(s)}
\]
\label{eq047}\[
  \mathrm{Cu(s)} + 2\,\mathrm{Ag^+(aq)} \longrightarrow \mathrm{Cu^{2+}(aq)} + 2\,\mathrm{Ag(s)}
\]
\label{eq048}\[
  \mathrm{Zn(s)} + 2\,\mathrm{Ag^+(aq)} \longrightarrow \mathrm{Zn^{2+}(aq)} + 2\,\mathrm{Ag(s)}
\]
Weak acids and bases ionise only partly, which is why equilibrium arrows are used:\nopagebreak
\label{eq049}\[
  \mathrm{HF(aq)} + \mathrm{H_2O(l)} \rightleftharpoons \mathrm{H_3O^+(aq)} + \mathrm{F^-(aq)}
\]
\label{eq050}\[
  \mathrm{HNO_2(aq)} + \mathrm{H_2O(l)} \rightleftharpoons \mathrm{H_3O^+(aq)} + \mathrm{NO_2^-(aq)}
\]
\label{eq051}\[
  \mathrm{CH_3COOH(aq)} + \mathrm{H_2O(l)} \rightleftharpoons \mathrm{H_3O^+(aq)} + \mathrm{CH_3COO^-(aq)}
\]
\label{eq052}\[
  \mathrm{HCN(aq)} + \mathrm{H_2O(l)} \rightleftharpoons \mathrm{H_3O^+(aq)} + \mathrm{CN^-(aq)}
\]
\label{eq053}\[
  \mathrm{HClO(aq)} + \mathrm{H_2O(l)} \rightleftharpoons \mathrm{H_3O^+(aq)} + \mathrm{ClO^-(aq)}
\]
\label{eq054}\[
  \mathrm{NH_3(aq)} + \mathrm{H_2O(l)} \rightleftharpoons \mathrm{NH_4^+(aq)} + \mathrm{OH^-(aq)}
\]
\label{eq055}\[
  \mathrm{CH_5N(aq)} + \mathrm{H_2O(l)} \rightleftharpoons \mathrm{CH_6N^+(aq)} + \mathrm{OH^-(aq)}
\]
Half-equations show the electrons that are gained or lost:\nopagebreak
\label{eq056}\[
  \mathrm{Cu^{2+}(aq)} + 2\,\mathrm{e^-} \longrightarrow \mathrm{Cu(s)}
\]
\label{eq057}\[
  \mathrm{Fe^{3+}(aq)} + \mathrm{e^-} \longrightarrow \mathrm{Fe^{2+}(aq)}
\]
\label{eq058}\[
  \mathrm{MnO_4^-(aq)} + 8\,\mathrm{H^+(aq)} + 5\,\mathrm{e^-} \longrightarrow \mathrm{Mn^{2+}(aq)} + 4\,\mathrm{H_2O(l)}
\]
\label{eq059}\[
  \mathrm{Zn(s)} \longrightarrow \mathrm{Zn^{2+}(aq)} + 2\,\mathrm{e^-}
\]
\label{eq060}\[
  2\,\mathrm{Cl^-(aq)} \longrightarrow \mathrm{Cl_2(g)} + 2\,\mathrm{e^-}
\]
\label{eq061}\[
  \mathrm{Al^{3+}(aq)} + 3\,\mathrm{e^-} \longrightarrow \mathrm{Al(s)}
\]
\label{eq062}\[
  \mathrm{O_2(g)} + 4\,\mathrm{H^+(aq)} + 4\,\mathrm{e^-} \longrightarrow 2\,\mathrm{H_2O(l)}
\]
\label{eq063}\[
  \mathrm{Cr_2O_7^{2-}(aq)} + 14\,\mathrm{H^+(aq)} + 6\,\mathrm{e^-} \longrightarrow 2\,\mathrm{Cr^{3+}(aq)} + 7\,\mathrm{H_2O(l)}
\]
Heating can break a compound down into simpler substances:\nopagebreak
\label{eq064}\[
  \mathrm{CaCO_3(s)} \xrightarrow{\Delta} \mathrm{CaO(s)} + \mathrm{CO_2(g)}
\]
\label{eq065}\[
  2\,\mathrm{KClO_3(s)} \xrightarrow{\Delta} 2\,\mathrm{KCl(s)} + 3\,\mathrm{O_2(g)}
\]
\label{eq066}\[
  2\,\mathrm{H_2O_2(aq)} \longrightarrow 2\,\mathrm{H_2O(l)} + \mathrm{O_2(g)}
\]
\label{eq067}\[
  2\,\mathrm{NaHCO_3(s)} \xrightarrow{\Delta} \mathrm{Na_2CO_3(s)} + \mathrm{H_2O(g)} + \mathrm{CO_2(g)}
\]
Many gas-phase reactions reach a dynamic equilibrium instead of going to completion:\nopagebreak
\label{eq068}\[
  \mathrm{N_2(g)} + 3\,\mathrm{H_2(g)} \rightleftharpoons 2\,\mathrm{NH_3(g)}
\]
\label{eq069}\[
  \mathrm{N_2O_4(g)} \rightleftharpoons 2\,\mathrm{NO_2(g)}
\]
\label{eq070}\[
  2\,\mathrm{SO_2(g)} + \mathrm{O_2(g)} \rightleftharpoons 2\,\mathrm{SO_3(g)}
\]
\label{eq071}\[
  \mathrm{H_2(g)} + \mathrm{I_2(g)} \rightleftharpoons 2\,\mathrm{HI(g)}
\]
\section*{Nuclear Equations}
In nuclear notation the mass number is a superscript and the atomic number a subscript, both before the symbol.
The nuclide \label{eq072}${}^{14}_{6}\mathrm{C}$ has 6 protons and 8 neutrons. The nuclide \label{eq073}${}^{238}_{92}\mathrm{U}$ has 92 protons and 146 neutrons. The nuclide \label{eq074}${}^{60}_{27}\mathrm{Co}$ has 27 protons and 33 neutrons. The nuclide \label{eq075}${}^{131}_{53}\mathrm{I}$ has 53 protons and 78 neutrons. The nuclide \label{eq076}${}^{56}_{26}\mathrm{Fe}$ has 26 protons and 30 neutrons. The nuclide \label{eq077}${}^{16}_{8}\mathrm{O}$ has 8 protons and 8 neutrons.
Alpha decay of U-238:\nopagebreak
\label{eq078}\[
  {}^{238}_{92}\mathrm{U} \longrightarrow {}^{234}_{90}\mathrm{Th} + {}^{4}_{2}\mathrm{He}
\]
Alpha decay of Ra-226:\nopagebreak
\label{eq079}\[
  {}^{226}_{88}\mathrm{Ra} \longrightarrow {}^{222}_{86}\mathrm{Rn} + {}^{4}_{2}\mathrm{He}
\]
Alpha decay of Po-210:\nopagebreak
\label{eq080}\[
  {}^{210}_{84}\mathrm{Po} \longrightarrow {}^{206}_{82}\mathrm{Pb} + {}^{4}_{2}\mathrm{He}
\]
Alpha decay of Pu-239:\nopagebreak
\label{eq081}\[
  {}^{239}_{94}\mathrm{Pu} \longrightarrow {}^{235}_{92}\mathrm{U} + {}^{4}_{2}\mathrm{He}
\]
Alpha decay of Rn-222:\nopagebreak
\label{eq082}\[
  {}^{222}_{86}\mathrm{Rn} \longrightarrow {}^{218}_{84}\mathrm{Po} + {}^{4}_{2}\mathrm{He}
\]
Alpha decay of Th-232:\nopagebreak
\label{eq083}\[
  {}^{232}_{90}\mathrm{Th} \longrightarrow {}^{228}_{88}\mathrm{Ra} + {}^{4}_{2}\mathrm{He}
\]
Alpha decay of Am-241:\nopagebreak
\label{eq084}\[
  {}^{241}_{95}\mathrm{Am} \longrightarrow {}^{237}_{93}\mathrm{Np} + {}^{4}_{2}\mathrm{He}
\]
Beta decay of C-14:\nopagebreak
\label{eq085}\[
  {}^{14}_{6}\mathrm{C} \longrightarrow {}^{14}_{7}\mathrm{N} + {}^{0}_{-1}\mathrm{e} + \bar{\nu}_{e}
\]
Beta decay of H-3:\nopagebreak
\label{eq086}\[
  {}^{3}_{1}\mathrm{H} \longrightarrow {}^{3}_{2}\mathrm{He} + {}^{0}_{-1}\mathrm{e} + \bar{\nu}_{e}
\]
Beta decay of Co-60:\nopagebreak
\label{eq087}\[
  {}^{60}_{27}\mathrm{Co} \longrightarrow {}^{60}_{28}\mathrm{Ni} + {}^{0}_{-1}\mathrm{e} + \bar{\nu}_{e}
\]
Beta decay of Sr-90:\nopagebreak
\label{eq088}\[
  {}^{90}_{38}\mathrm{Sr} \longrightarrow {}^{90}_{39}\mathrm{Y} + {}^{0}_{-1}\mathrm{e} + \bar{\nu}_{e}
\]
Beta decay of I-131:\nopagebreak
\label{eq089}\[
  {}^{131}_{53}\mathrm{I} \longrightarrow {}^{131}_{54}\mathrm{Xe} + {}^{0}_{-1}\mathrm{e} + \bar{\nu}_{e}
\]
Beta decay of Cs-137:\nopagebreak
\label{eq090}\[
  {}^{137}_{55}\mathrm{Cs} \longrightarrow {}^{137}_{56}\mathrm{Ba} + {}^{0}_{-1}\mathrm{e} + \bar{\nu}_{e}
\]
Beta decay of P-32:\nopagebreak
\label{eq091}\[
  {}^{32}_{15}\mathrm{P} \longrightarrow {}^{32}_{16}\mathrm{S} + {}^{0}_{-1}\mathrm{e} + \bar{\nu}_{e}
\]
Beta decay of Pb-210:\nopagebreak
\label{eq092}\[
  {}^{210}_{82}\mathrm{Pb} \longrightarrow {}^{210}_{83}\mathrm{Bi} + {}^{0}_{-1}\mathrm{e} + \bar{\nu}_{e}
\]
Positron emission by C-11:\nopagebreak
\label{eq093}\[
  {}^{11}_{6}\mathrm{C} \longrightarrow {}^{11}_{5}\mathrm{B} + {}^{0}_{1}\mathrm{e} + \nu_{e}
\]
Positron emission by N-13:\nopagebreak
\label{eq094}\[
  {}^{13}_{7}\mathrm{N} \longrightarrow {}^{13}_{6}\mathrm{C} + {}^{0}_{1}\mathrm{e} + \nu_{e}
\]
Positron emission by O-15:\nopagebreak
\label{eq095}\[
  {}^{15}_{8}\mathrm{O} \longrightarrow {}^{15}_{7}\mathrm{N} + {}^{0}_{1}\mathrm{e} + \nu_{e}
\]
Positron emission by F-18:\nopagebreak
\label{eq096}\[
  {}^{18}_{9}\mathrm{F} \longrightarrow {}^{18}_{8}\mathrm{O} + {}^{0}_{1}\mathrm{e} + \nu_{e}
\]
Positron emission by Na-22:\nopagebreak
\label{eq097}\[
  {}^{22}_{11}\mathrm{Na} \longrightarrow {}^{22}_{10}\mathrm{Ne} + {}^{0}_{1}\mathrm{e} + \nu_{e}
\]
Electron capture in Be-7:\nopagebreak
\label{eq098}\[
  {}^{7}_{4}\mathrm{Be} + {}^{0}_{-1}\mathrm{e} \longrightarrow {}^{7}_{3}\mathrm{Li} + \nu_{e}
\]
Electron capture in Ar-37:\nopagebreak
\label{eq099}\[
  {}^{37}_{18}\mathrm{Ar} + {}^{0}_{-1}\mathrm{e} \longrightarrow {}^{37}_{17}\mathrm{Cl} + \nu_{e}
\]
Electron capture in Fe-55:\nopagebreak
\label{eq100}\[
  {}^{55}_{26}\mathrm{Fe} + {}^{0}_{-1}\mathrm{e} \longrightarrow {}^{55}_{25}\mathrm{Mn} + \nu_{e}
\]
Electron capture in Cr-51:\nopagebreak
\label{eq101}\[
  {}^{51}_{24}\mathrm{Cr} + {}^{0}_{-1}\mathrm{e} \longrightarrow {}^{51}_{23}\mathrm{V} + \nu_{e}
\]
Gamma emission from an excited Ni-60 nucleus:\nopagebreak
\label{eq102}\[
  {}^{60}_{28}\mathrm{Ni}^{*} \longrightarrow {}^{60}_{28}\mathrm{Ni} + \gamma
\]
Gamma emission from an excited Sr-87 nucleus:\nopagebreak
\label{eq103}\[
  {}^{87}_{38}\mathrm{Sr}^{*} \longrightarrow {}^{87}_{38}\mathrm{Sr} + \gamma
\]
Fission of uranium-235 after neutron capture:\nopagebreak
\label{eq104}\[
  {}^{235}_{92}\mathrm{U} + {}^{1}_{0}\mathrm{n} \longrightarrow {}^{141}_{56}\mathrm{Ba} + {}^{92}_{36}\mathrm{Kr} + 3\,{}^{1}_{0}\mathrm{n}
\]
Deuterium and tritium fuse:\nopagebreak
\label{eq105}\[
  {}^{2}_{1}\mathrm{H} + {}^{3}_{1}\mathrm{H} \longrightarrow {}^{4}_{2}\mathrm{He} + {}^{1}_{0}\mathrm{n}
\]
Nitrogen bombarded with alpha particles gives oxygen:\nopagebreak
\label{eq106}\[
  {}^{14}_{7}\mathrm{N} + {}^{4}_{2}\mathrm{He} \longrightarrow {}^{17}_{8}\mathrm{O} + {}^{1}_{1}\mathrm{H}
\]
The first artificial radioactive isotope was made from aluminium:\nopagebreak
\label{eq107}\[
  {}^{27}_{13}\mathrm{Al} + {}^{4}_{2}\mathrm{He} \longrightarrow {}^{30}_{15}\mathrm{P} + {}^{1}_{0}\mathrm{n}
\]
\section*{Measurements and Units}
Quantities in the laboratory carry units, and many of those units have powers.
At sea level the density of air is about 1.2 kg m$^{-3}$, and a barometer reading corresponds to a pressure of 1.1 times 10\textsuperscript{5} N/m\textsuperscript{2}. The solution contains 0.1 mol dm$^{-3}$ of hydrochloric acid, so 250 cm\textsuperscript{3} of it holds 0.025 mol. The reservoir holds 2.5 km\textsuperscript{3} of water spread over 150 km\textsuperscript{2}. Near the surface of the earth, free fall has an acceleration of 9.81 m s$^{-2}$.

The intensity of sunlight at the ground is roughly 950 W m$^{-2}$ on a clear day. The density of ice is 0.92 g/cm\textsuperscript{3}, so a block of volume 60 cm\textsuperscript{3} has a mass of 55 g. The pump delivers 0.8 m\textsuperscript{3}/s, which is 800 dm\textsuperscript{3} every second. The specific heat capacity of aluminium is 0.9 J g$^{-1}$ K$^{-1}$.

The plate has an area of 0.5 m\textsuperscript{2} and a thickness of 10 cm, so its volume is 0.05 m\textsuperscript{3}. A steel wire with a cross-section of 10.0 mm\textsuperscript{2} carries a load of 900 N, which is a stress of 90 N/mm\textsuperscript{2}. The measured viscosity of the ice melt is 0.5 kg m$^{-1}$ s$^{-1}$. A cart accelerates at 4.5 m/s\textsuperscript{2} for 6 s and reaches a speed of 27 m/s.

The plate has an area of 2.0 m\textsuperscript{2} and a thickness of 2 cm, so its volume is 0.04 m\textsuperscript{3}. The pump delivers 0.4 m\textsuperscript{3}/s, which is 400 dm\textsuperscript{3} every second. The measured viscosity of the aluminium melt is 0.5 kg m$^{-1}$ s$^{-1}$. The specific heat capacity of aluminium is 0.9 J g$^{-1}$ K$^{-1}$.

A cart accelerates at 3.0 m/s\textsuperscript{2} for 4 s and reaches a speed of 12 m/s. The solution contains 2.0 mol dm$^{-3}$ of sodium chloride, so 250 cm\textsuperscript{3} of it holds 0.5 mol. Near the surface of the earth, free fall has an acceleration of 9.81 m s$^{-2}$. The density of ice is 0.92 g/cm\textsuperscript{3}, so a block of volume 60 cm\textsuperscript{3} has a mass of 55 g.

The intensity of sunlight at the ground is roughly 1020 W m$^{-2}$ on a clear day. The reservoir holds 6.0 km\textsuperscript{3} of water spread over 150 km\textsuperscript{2}. A steel wire with a cross-section of 2.0 mm\textsuperscript{2} carries a load of 900 N, which is a stress of 450 N/mm\textsuperscript{2}. At sea level the density of air is about 1.2 kg m$^{-3}$, and a barometer reading corresponds to a pressure of 2.5 times 10\textsuperscript{5} N/m\textsuperscript{2}.

The specific heat capacity of copper is 0.39 J g$^{-1}$ K$^{-1}$. At sea level the density of air is about 1.2 kg m$^{-3}$, and a barometer reading corresponds to a pressure of 1.2 times 10\textsuperscript{5} N/m\textsuperscript{2}. Near the surface of the earth, free fall has an acceleration of 9.8 m s$^{-2}$. The solution contains 0.1 mol dm$^{-3}$ of sodium chloride, so 250 cm\textsuperscript{3} of it holds 0.025 mol.

A steel wire with a cross-section of 8.0 mm\textsuperscript{2} carries a load of 900 N, which is a stress of 112.5 N/mm\textsuperscript{2}. A cart accelerates at 2.0 m/s\textsuperscript{2} for 6 s and reaches a speed of 12 m/s.
\section*{Atoms and Nuclei in Running Text}
The mass number of a nuclide stands before its symbol, as a raised number.
The nucleus \textsuperscript{14}C decays by beta emission, and its half-life \textit{t}$_{1/2}$ is 5730 years. Cobalt-60, written \textsuperscript{60}Co, contains 27 protons and 33 neutrons. The isotopes \textsuperscript{131}I and \textsuperscript{129}I of iodine have the same chemical behaviour but different masses. A sample of \textsuperscript{137}Cs loses half of its atoms every 30.2 years, so the activity falls in a regular pattern.

The nucleus \textsuperscript{226}Ra decays by alpha emission, and its half-life \textit{t}$_{1/2}$ is 1600 years. Polonium-210, written \textsuperscript{210}Po, contains 84 protons and 126 neutrons. The isotopes \textsuperscript{238}U and \textsuperscript{239}U of uranium have the same chemical behaviour but different masses. A sample of \textsuperscript{241}Am loses half of its atoms every 432 years, so the activity falls in a regular pattern.

The nucleus \textsuperscript{32}P decays by beta emission, and its half-life \textit{t}$_{1/2}$ is 14.3 days. Hydrogen-3, written \textsuperscript{3}H, contains 1 proton and 2 neutrons. The isotopes \textsuperscript{222}Rn and \textsuperscript{221}Rn of radon have the same chemical behaviour but different masses. A sample of \textsuperscript{90}Sr loses half of its atoms every 28.8 years, so the activity falls in a regular pattern.

The nucleus \textsuperscript{14}C decays by beta emission, and its half-life \textit{t}$_{1/2}$ is 5730 years. Cobalt-60, written \textsuperscript{60}Co, contains 27 protons and 33 neutrons. The isotopes \textsuperscript{131}I and \textsuperscript{130}I of iodine have the same chemical behaviour but different masses. A sample of \textsuperscript{137}Cs loses half of its atoms every 30.2 years, so the activity falls in a regular pattern.

The nucleus \textsuperscript{226}Ra decays by alpha emission, and its half-life \textit{t}$_{1/2}$ is 1600 years. Polonium-210, written \textsuperscript{210}Po, contains 84 protons and 126 neutrons. The isotopes \textsuperscript{238}U and \textsuperscript{236}U of uranium have the same chemical behaviour but different masses. A sample of \textsuperscript{241}Am loses half of its atoms every 432 years, so the activity falls in a regular pattern.

The nucleus \textsuperscript{32}P decays by beta emission, and its half-life \textit{t}$_{1/2}$ is 14.3 days. Hydrogen-3, written \textsuperscript{3}H, contains 1 proton and 2 neutrons. The isotopes \textsuperscript{222}Rn and \textsuperscript{223}Rn of radon have the same chemical behaviour but different masses. A sample of \textsuperscript{90}Sr loses half of its atoms every 28.8 years, so the activity falls in a regular pattern.
\section*{Other Scripts in Running Text}
Indices, powers of ten and ordinal endings also appear inside sentences.
The kinetic energy \textit{E}$_{k}$ of the cart is 6 J at the bottom of the ramp. Expanding (\textit{a} + \textit{b})\textsuperscript{2} gives \textit{a}\textsuperscript{2} + 2\textit{ab} + \textit{b}\textsuperscript{2}. A tolerance of 10$^{-3}$ is enough for the 12\textsuperscript{th} iteration of the method. The roots of the quadratic are \textit{x}\textsubscript{1} = 2 and \textit{x}\textsubscript{2} = 5.

The coefficient of determination \textit{R}\textsuperscript{2} equals 0.84 for this fit. Each term satisfies \textit{a}$_{n+1}$ = \textit{a}$_{n}$ + 7, so with \textit{a}\textsubscript{1} = 2 the 5\textsuperscript{th} term is 30. The null hypothesis \textit{H}\textsubscript{0} is rejected when the p-value is below 0.05. The 1\textsuperscript{st}, 2\textsuperscript{nd} and 3\textsuperscript{rd} terms are 5, 9 and 7.

The point (\textit{x}\textsubscript{1}, \textit{y}\textsubscript{1}) lies on the line, and so does (\textit{x}\textsubscript{2}, \textit{y}\textsubscript{2}). The square of the standard deviation, \textit{s}\textsuperscript{2}, is called the variance. The half-life \textit{t}$_{1/2}$ of the sample is 8 hours. Taking logarithms to base ten, log\textsubscript{10} 3000 is about 3.48.

The mass of the sample is 1 times 10$^{-3}$ kg. The initial velocity \textit{v}\textsubscript{0} is 7 m/s and the final velocity \textit{v}$_{f}$ is 10 m/s. The coefficient of determination \textit{R}\textsuperscript{2} equals 0.91 for this fit. The kinetic energy \textit{E}$_{k}$ of the cart is 15 J at the bottom of the ramp.

Each term satisfies \textit{a}$_{n+1}$ = \textit{a}$_{n}$ + 5, so with \textit{a}\textsubscript{1} = 3 the 5\textsuperscript{th} term is 23. The null hypothesis \textit{H}\textsubscript{0} is rejected when the p-value is below 0.05. The roots of the quadratic are \textit{x}\textsubscript{1} = 5 and \textit{x}\textsubscript{2} = 3. The half-life \textit{t}$_{1/2}$ of the sample is 8 hours.

The square of the standard deviation, \textit{s}\textsuperscript{2}, is called the variance. The 1\textsuperscript{st}, 2\textsuperscript{nd} and 3\textsuperscript{rd} terms are 7, 2 and 9. A tolerance of 10$^{-3}$ is enough for the 17\textsuperscript{th} iteration of the method. The point (\textit{x}\textsubscript{1}, \textit{y}\textsubscript{1}) lies on the line, and so does (\textit{x}\textsubscript{2}, \textit{y}\textsubscript{2}).

Expanding (\textit{a} + \textit{b})\textsuperscript{2} gives \textit{a}\textsuperscript{2} + 2\textit{ab} + \textit{b}\textsuperscript{2}. The initial velocity \textit{v}\textsubscript{0} is 7 m/s and the final velocity \textit{v}$_{f}$ is 10 m/s. The mass of the sample is 1.5 times 10$^{-3}$ kg. Taking logarithms to base ten, log\textsubscript{10} 3000 is about 3.48.

Each term satisfies \textit{a}$_{n+1}$ = \textit{a}$_{n}$ + 3, so with \textit{a}\textsubscript{1} = 7 the 5\textsuperscript{th} term is 19. The kinetic energy \textit{E}$_{k}$ of the cart is 21 J at the bottom of the ramp. The 1\textsuperscript{st}, 2\textsuperscript{nd} and 3\textsuperscript{rd} terms are 7, 5 and 3. The coefficient of determination \textit{R}\textsuperscript{2} equals 0.91 for this fit.

Taking logarithms to base ten, log\textsubscript{10} 200 is about 2.30. The point (\textit{x}\textsubscript{1}, \textit{y}\textsubscript{1}) lies on the line, and so does (\textit{x}\textsubscript{2}, \textit{y}\textsubscript{2}).
\end{document}
